%% Elementary PCF proofs; the maximum theorem is a separately cited result.

\section{Products and possible cofinalities}

\begin{annexitem}{annex:p6.pcf.1}{thm:p06-pcf-elementary}{Elementary properties of possible cofinalities}
\begin{proof}
\paragraph{The quotient order.}
First let \(D\) be a proper filter.
Equality modulo \(D\) is reflexive and symmetric.
If \(f={}_D g\) and \(g={}_D h\), the intersection of their two
agreement sets belongs to \(D\) and is contained in the agreement
set of \(f,h\).  Closure under supersets gives \(f={}_D h\).
Thus it is an equivalence relation.
Comparison modulo \(D\) is reflexive, and the same intersection
argument proves transitivity.
If \(f={}_D f'\) and \(g={}_D g'\), intersect the two agreement
sets with \(\{\kappa:f(\kappa)\leq g(\kappa)\}\).
It follows that \(f\leq_D g\) implies \(f'\leq_D g'\), and the
reverse follows symmetrically.
The quotient comparison is therefore well defined.
If both \(f\leq_D g\) and \(g\leq_D f\), their comparison sets
intersect inside the equality set, so \(f={}_D g\).
This proves antisymmetry on the quotient.

Now assume \(D\) is an ultrafilter.
The three sets where \(f(\kappa)<g(\kappa)\),
\(f(\kappa)=g(\kappa)\), and \(f(\kappa)>g(\kappa)\) partition
\(\mathfrak a\).
At most one can belong to a proper filter, because they are disjoint.
At least one belongs to \(D\): otherwise all three complements
would belong to \(D\) and their intersection would be empty.
This gives linearity and identifies \(<_D\) with strict comparison
between quotient classes.

The zero function belongs to the product, so the quotient is nonempty.
For any \(f\in\prod\mathfrak a\), set
\[
 g(\kappa)=f(\kappa)+1.
\]
Every infinite cardinal is a limit ordinal, so \(g(\kappa)<\kappa\).
Thus \(g\) also belongs to the product, and \(f<_D g\).
The quotient has no greatest element.

\paragraph{Cofinality of a linear order without a greatest element.}
Let \(L\) be any nonempty linear order without a greatest element,
and let \(\lambda\) be the least cardinality of a cofinal subset.
This cardinal is infinite: a finite subset has a largest member,
and an element strictly above that member is not dominated by it.
Choose a cofinal family \(\{c_\xi:\xi<\lambda\}\).

Recursively construct a strictly increasing sequence
\(\langle d_\xi:\xi<\lambda\rangle\) such that \(c_\xi<d_\xi\).
At stage \(\xi\), the set
\[
 \{d_\eta:\eta<\xi\}\cup\{c_\xi\}
\]
has cardinality below \(\lambda\), because \(\lambda\) is an infinite
cardinal and \(\xi<\lambda\).
It is not cofinal by minimality of \(\lambda\).
In a linear order this means that some element is strictly above
every member of that set; choose it as \(d_\xi\).
The recursion is valid without assuming regularity of \(\lambda\).
Its resulting sequence is cofinal since it dominates the chosen
cofinal family.

If \(\lambda\) were singular, a cofinal subset \(I\subseteq\lambda\)
of cardinality \(\operatorname{cf}(\lambda)<\lambda\) would make
\(\{d_\xi:\xi\in I\}\) cofinal in \(L\).
Indeed, after any sequence index there is one in \(I\), and the
sequence is increasing.  This contradicts minimality of \(\lambda\).
Hence \(\lambda\) is regular.
Apply this argument to \(\prod\mathfrak a/D\), and choose one
representative for each element of the resulting sequence.
The representatives form a \(<_D\)-increasing, \(\leq_D\)-cofinal
sequence of functions of length \(\lambda\).

\paragraph{Principal and finite cases.}
For \(\kappa\in\mathfrak a\), let
\[
 D_\kappa=\{b\subseteq\mathfrak a:\kappa\in b\}.
\]
This is a proper ultrafilter.
The map
\[
 \prod\mathfrak a/D_\kappa\longrightarrow\kappa,
 \qquad [f]_{D_\kappa}\longmapsto f(\kappa)
\]
is well defined, injective, and preserves and reflects order.
It is onto: for any \(\beta<\kappa\), use the function with value
\(\beta\) at \(\kappa\) and zero at every other coordinate.
Thus the quotient is order-isomorphic to \(\kappa\), whose
cofinality is \(\kappa\) by regularity.
It follows that \(\mathfrak a\subseteq\operatorname{pcf}(\mathfrak a)\).

If \(\mathfrak a\) is finite, an ultrafilter on it must contain a
singleton.  Otherwise it contains the complement of every singleton,
whose finite intersection is empty.
If \(\{\kappa\}\in D\), closure under supersets gives
\(D_\kappa\subseteq D\), and properness forbids any member of \(D\)
missing \(\kappa\).  Hence \(D=D_\kappa\).
All its possible cofinalities are therefore exactly the members of
\(\mathfrak a\).
\end{proof}
\end{annexitem}

\begin{annexitem}{annex:p6.pcf.2}{thm:p06-pcf-bounds}{Bounds at a singular limit of regular cardinals}
\begin{proof}
Write \(\lambda=\max\operatorname{pcf}(\mathfrak a)\), whose existence
is supplied by the cited maximum theorem,
Theorem~\ref{thm:p06-pcf-maximum}.
It equals the least size of a pointwise cofinal family in
\(\prod\mathfrak a\).
If \(F\) is any such family, its projection to each coordinate
\(\kappa\in\mathfrak a\) is cofinal in \(\kappa\).
For each \(\beta<\kappa\), the product contains a function having
value \(\beta\) at that coordinate and zero elsewhere, and some
member of \(F\) dominates that function.
Regularity of \(\kappa\) therefore gives \(|F|\geq\kappa\).
Taking all coordinates yields \(|F|\geq\sup\mathfrak a\).
The whole product is itself cofinal, so
\[
 \sup\mathfrak a\leq\lambda\leq\left|\prod\mathfrak a\right|.
\]

Suppose that \(\mathfrak a\) has no largest element, and put
\(\mu=\sup\mathfrak a\).
The set \(\mathfrak a\) is cofinal in the cardinal \(\mu\), so
\[
 \operatorname{cf}(\mu)\leq|\mathfrak a|
             <\min\mathfrak a\leq\mu.
\]
Thus \(\mu\) is singular.
By Lemma~\ref{thm:p06-pcf-elementary}, every member of
\(\operatorname{pcf}(\mathfrak a)\), including its maximum
\(\lambda\), is regular.
The inequality \(\lambda\geq\mu\) therefore strengthens to
\(\lambda>\mu\).

Now take \(\mathfrak a=\{\aleph_{n+1}:n<\omega\}\).
Its factors are regular successor cardinals, its size is \(\aleph_0\),
and its least member is \(\aleph_1\), so it is progressive.
It has no largest element and has supremum \(\aleph_\omega\).
The strict lower bound gives
\(\lambda\geq\aleph_{\omega+1}\).
Enumerating the coordinates identifies the product with a subset
of the functions from \(\omega\) into \(\aleph_\omega\); hence
\[
 \lambda\leq\left|\prod\mathfrak a\right|
       \leq\aleph_\omega^{\aleph_0}
       \leq(2^{\aleph_\omega})^{\aleph_0}
       =2^{\aleph_\omega}.
\]
The last equality is the cardinal exponent law and
\(\aleph_\omega\cdot\aleph_0=\aleph_\omega\).
If \(2^{\aleph_\omega}=\aleph_{\omega+1}\), the lower and upper
bounds coincide, proving the final assertion.
\end{proof}
\end{annexitem}
